\scanpage{27}

\uline{Cor. 2} \unsure{Èquiv.|Équiv.} de cat. fibrées entre

$1^\circ$) $(E,Q)$, $Q$ une forme quadratique binaire qui n'est nulle sur aucune fibre réduite

$2^\circ$) Systèmes $(\mathcal A,\Pi)$, où $\mathcal A$ est une alg. quadratique sur $S$, et où $\Pi$ est un torseur\textsuperscript{*} sous le groupe unitaire $U_{\mathcal A}$. \deleted{\unsure{*}}

% slanted
\margin{Torseur \unsure{sens|sous} fppf ou \unsure{sens|sous} (ét), cela revient ici au $=$, mais la top. (Zar) n'est plus assez fine \ldots}

$\Pi$ est le sous-schéma de $E^{\unsure{*}}$ des sections locales $x$ telles que $Q(x)=1$, qui \deleted{sont}\inserted{forment bien} un torseur sous $U_{\mathcal A}$.

\uline{NB} Le rev. quadratique associé à la forme quadratique $Q:E\longrightarrow\Delta'$ (non nulle sur toute fibre) peut s'expliciter comme \inserted{le} sous-schéma de ${\unsure{P|\mathbb P}}^1(E)$ défini par l'équation $Q(x)=0$. Si $E$ est considéré comme Module inv. sur l'Algèbre quadratique $\mathcal A$, le morphisme

\[
\operatorname{Spec}(\mathcal A)\lhook\joinrel\longrightarrow{P}^{\unsure{*1|1}}(E)
\]

correspondant \deleted{peut} s'explicite ainsi : pour toute section locale $u$ de $\operatorname{Spec}(\mathcal A)$ sur $A$ i.e.
% archive sheet number: 26
